\section{Introduction}
\label{sec:introduction}

Discrete element method \citep[DEM;][]{Cundall1979} simulations have become an indispensable tool for investigating the micro-mechanical origins of soil behavior, including liquefaction resistance \citep{Huang2018,Yang2021,WuYang2022}, fabric evolution \citep{Zhao2024,YangTaiebat2024}, and multi-directional cyclic loading effects \citep{Wei2020,Yang2022}. As DEM studies grow in scope and complexity, the associated workflow---writing simulation scripts, managing long-running jobs, querying model state, and iterating on parameters---has become a significant bottleneck that limits research productivity. Recent advances in large language model (LLM) agents offer a potential path to alleviating this burden, yet existing integrations between LLMs and engineering simulation software remain limited in scope.

The application of LLMs to geotechnical engineering has progressed rapidly. \citet{Xu2024} proposed a specialized LLM framework for pile design that combined prompt engineering with an external mathematical computation module. \citet{Ding2025} extended this to a multimodal multi-agent system capable of processing both text and image inputs for footing design. \citet{Kim2024} explored ChatGPT-assisted MATLAB code generation for seepage and slope stability analyses. In the broader computational mechanics community, \citet{GeoSimAI2025} presented GeoSim.AI, an LLM assistant with retrieval-augmented generation (RAG) for translating natural language instructions into geomechanical simulation scripts. Beyond geotechnical applications, agent-based frameworks have been proposed for molecular dynamics simulations \citep{AtomAgents2024,LangSim2024}, computational fluid dynamics \citep{FoamAgent2025}, and density functional theory \citep{VASPilot2025}.

Despite this progress, three fundamental limitations persist across existing approaches:
\begin{enumerate}
  \item \textbf{One-directional communication.} Existing tools generate commands or scripts that the user executes manually. The LLM has no mechanism to observe the simulation state during or after execution, preventing iterative diagnosis and parameter adjustment.
  \item \textbf{No runtime integration.} Connection to live simulation processes is either absent or described as future work. GeoSim.AI explicitly states that ``connection of the LLM to the repository, for example to run simulation using a chosen software or extract results from an already ran simulation, is under test and development'' \citep{GeoSimAI2025}.
  \item \textbf{Duplicated orchestration effort.} Several projects build custom agent orchestration layers---conversation management, tool routing, multi-provider LLM support---that are rapidly superseded by general-purpose frameworks from model providers, diverting development effort away from domain-specific contributions.
\end{enumerate}

The present study introduces \texttt{pfc-mcp}, an open-source Model Context Protocol \citep[MCP;][]{MCP2024} server that provides real-time bidirectional integration between LLM agents and ITASCA PFC \citep{PFC2023}, a widely used commercial DEM software. Rather than building a standalone agent platform, \texttt{pfc-mcp} exposes domain-specific capabilities as MCP tools that any MCP-compatible LLM client can consume, delegating all orchestration to the upstream platform. The server comprises two components: a documentation subsystem providing structured access to PFC commands, Python API, and reference documentation; and an execution subsystem backed by a bridge process running inside the PFC GUI, enabling both synchronous REPL-style queries and asynchronous long-running simulation tasks. Critically, the REPL interface remains available whether PFC is actively cycling or idle, enabling live monitoring of ongoing simulations without interruption.

The system design and its applicability are demonstrated through a complete liquefaction simulation workflow in which an LLM agent diagnoses a numerical instability in a cyclic shear simulation, queries the PFC documentation to identify the relevant command, modifies the simulation script, submits the corrected job, and monitors pore pressure generation and contact evolution in real time. To the best of the authors' knowledge, \texttt{pfc-mcp} is the first MCP server for any geotechnical or DEM simulation software, and the first system to enable real-time bidirectional LLM--DEM interaction.

The remainder of this paper is organized as follows. \Cref{sec:design} presents the system architecture, the bridge runtime scheduling mechanism, and the design rationale behind key decisions. \Cref{sec:demonstration} describes the liquefaction simulation workflow used to validate the system and discusses the agent--simulation interaction patterns observed. \Cref{sec:conclusions} summarizes the main contributions and identifies directions for future development.
